\section{犯罪形态 - 犯罪预备、未遂}
\begin{enumerate}
  \item 犯罪在时间上的发展问题。仅针对故意犯罪，过失只有成立、不成立（\bs{产生损害后果是过失犯罪的构成要件}）。
  \item 从阶段分。
        \begin{enumerate}
          \item 预备阶段 \yjt 预备、\zd{（预备）中止 \yjt 免除处罚}。
          \item 实行阶段 \yjt（实行）中止、未遂、既遂。
        \end{enumerate}
  \item 从完成形态分。
        \begin{enumerate}
          \item 未完成形态 \yjt 预备、中止（预备阶段 \hz 实行阶段）、未遂。
          \item 完成形态 \yjt 既遂。
        \end{enumerate}
  \item 预备。已开始未着手，意志以外的原因。
        \begin{enumerate}
          \item 预备行为。制造条件。有那么一丢丢对法益侵犯的风险。区分生活行为 + 犯意表示。
          \item 着手。现实（客观存在）、紧迫、直接危险。预备阶段 \yjt 实行阶段。开始 \bdy 着手，考时间差。
          \item 隔离犯、间接正犯、不作为犯 \yjt 具体分析。
        \end{enumerate}
  \item 未遂。不能继续而放弃。
        \begin{enumerate}
          \item \zd{vs 不能犯}。对法益没有危险。其实就是判断是不是一个危害行为。
                \begin{enumerate}
                  \item 看客观不看主观。不然迷信犯。\al{甲将\xhx{面粉}交由乙运输，谎称是毒品，乙以为是毒品就去运输}。
                  \item \zd{行为时一般人}。看彼时彼刻，不能看此时此刻。\al{投毒 1 毫克得投 1000000 天 \yjt 事后诸葛亮}。
                  \item \zd{少数说。事后科学角度判断 \yjt 手段不能犯}。
                  \item 辩证的角度。全面、联系、发展的眼光，不能片面。\al{手枪\xhx{射程 90 米}，朝 \xhx{100 米外}仇人开枪}。
                \end{enumerate}
        \end{enumerate}
\end{enumerate}
